%% Jian Zhao - Two-page master resume
%% Compressed from main_example-jian.tex; original preserved unchanged.
%% Compile with: lualatex -interaction=nonstopmode -halt-on-error main_example-jian-2page.tex

\documentclass[11pt,a4paper,sans]{moderncv}
\moderncvstyle{banking}
\moderncvcolor{blue}

\renewcommand*{\firstnamestyle}[1]{{\fontsize{34}{36}\bfseries\upshape\color{color1}#1}}
\renewcommand*{\lastnamestyle}[1]{{\fontsize{34}{36}\bfseries\upshape\color{color1}#1}}
\renewcommand*{\sectionstyle}[1]{{\sectionfont\color{color1}#1}}

\usepackage[utf8]{inputenc}
\usepackage{hyperref}
\hypersetup{
  colorlinks=true,
  urlcolor=blue,
  pdftitle={Jian Zhao - Resume},
  pdfpagemode=UseNone,
}
\usepackage[scale=0.80]{geometry}
\usepackage{needspace}

\name{Jian}{Zhao}
\address{Philadelphia, PA}{}{}
\phone[mobile]{302-359-7657}
\email{zhaojianphd@gmail.com}
\extrainfo{\href{https://linkedin.com/in/jian-zhao-12800616a}{LinkedIn}}

\begin{document}
\makecvtitle

\vspace{-2pt}
\small{AI researcher and software engineer with Tsinghua University graduate training, U.S. master's and Ph.D. research, and 10+ years spanning biomedical AI, graph learning, computer vision, LLM applications, HPC, and enterprise systems. Combines interdisciplinary healthcare research with production software engineering, multi-institution collaboration, and hands-on Python and PyTorch development.}

\section{Core Competencies}
\vspace{-1pt}
\begin{itemize}
\item \textbf{Generative AI \& Healthcare NLP}: LLaMA, ChatGPT API, Hugging Face Transformers, PEFT (LoRA/QLoRA), RAG, prompt engineering; Milvus, Pinecone, Neo4j; EMMA postpartum-depression support assistant
\item \textbf{Deep Learning \& Biomedical Vision}: Python, PyTorch, TensorFlow, Scikit-learn, PyTorch Geometric; CNNs, GNNs, representation learning; microscopy cell segmentation, tracking, and quantification
\item \textbf{Machine Learning \& Data}: XGBoost, Random Forest, SVM; clinical and public-health analytics, EHR workflows; SQL, PostgreSQL, Oracle, MySQL; Git, Linux, Docker
\item \textbf{HPC \& Engineering}: MPI, MPI I/O, OpenMP, GPU computing, OSIRIS, GROMACS, LUSTRE, Hadoop; Java, C++, C\#, MATLAB, R; enterprise software architecture and delivery
\end{itemize}

\section{Professional and Research Experience}
\vspace{1pt}
\begin{itemize}

\needspace{5\baselineskip}
\item{\cventry{Oct 2023--Apr 2025}{Assistant Research Scientist}{Temple University, College of Public Health}{Philadelphia, PA}{}{\vspace{1pt}
\begin{itemize}
  \item Developed EMMA, an LLM-based assistant for postpartum-depression support using LLaMA/ChatGPT, RAG, PEFT, Milvus, Neo4j, Flowise, Gradio, and Streamlit to connect users with locally relevant resources
  \item Applied Random Forest, XGBoost, and SVM to ambulatory-visit and public-health data, identifying blood-pressure and A1C monitoring gaps and supporting hypertension-prediction research
  \item Designed AI and data workflows spanning EHR, environmental sensors, and public-health repositories; contributed manuscripts, grant proposals, technical documentation, and cross-disciplinary research supervision
\end{itemize}}}

\vspace{2pt}
\needspace{5\baselineskip}
\item{\cventry{Aug 2019--Present}{Ph.D. Research Assistant}{Delaware State University}{Dover, DE}{}{\vspace{1pt}
\begin{itemize}
  \item Research biomedical cell segmentation, tracking, and quantification using mathematical algorithms, image processing, CNNs, GNNs, and LLM-assisted parameter tuning for microscopy applications
  \item Developed a rapid-tuning approach for CSTQ parameters supporting reproducible biomedical image-analysis and disease-research workflows
  \item Conduct large-scale nano-bio spectroscopy and plasma simulations using MPI, OpenMP, GPU computing, OSIRIS, and GROMACS in collaborations with Princeton, UCLA, and Lawrence Livermore National Laboratory
\end{itemize}}}

\vspace{2pt}
\needspace{5\baselineskip}
\item{\cventry{Aug 2016--Jul 2019}{Graduate Research Assistant}{Delaware State University}{Dover, DE}{}{\vspace{1pt}
\begin{itemize}
  \item Built and administered an HPC research cluster with 8 compute nodes, 320 CPUs, and 512 GB RAM for nano-bio spectroscopy simulations and data-intensive experiments
  \item Completed ML projects in pattern recognition and dimensionality reduction (PCA, LDA, LLE), and investigated visible-light access control for wireless insulin-pump security
\end{itemize}}}

\vspace{2pt}
\needspace{5\baselineskip}
\item{\cventry{Aug 2007--Jul 2016}{Senior Software Engineer}{PetroChina Planning \& Engineering Institute}{Beijing, China}{}{\vspace{1pt}
\begin{itemize}
  \item Developed BP neural-network and decision-tree models for natural-gas pipeline energy forecasting, achieving approximately 90\% accuracy for strategic planning analyses
  \item Built an M/M/C vehicle-allocation model and delivered analytics and enterprise software using MATLAB, Java, SQL, Weka, and statistical modeling
\end{itemize}}}

\vspace{2pt}
\needspace{5\baselineskip}
\item{\cventry{Aug 2005--Jul 2007}{Project Manager}{Chinese Academy of Sciences}{Beijing, China}{}{\vspace{1pt}
\begin{itemize}
  \item Led a five-person Sino-German national forest-pest-control information-system project, delivering a scalable browser/server application for the National Forestry program
  \item Modernized Forest News from ASP to JSP and migrated Access to Oracle 9i using Java, Spring, Hibernate, Struts, JUnit, and GIS technologies
\end{itemize}}}

\end{itemize}


\section{Education}
\vspace{0pt}
\begin{itemize}
\needspace{4\baselineskip}
\item{\cventry{Aug 2019--May 2027}{Ph.D. Candidate, Applied Mathematics and Mathematical Physics}{Delaware State University}{Dover, DE}{}{GPA: 3.89; research in biomedical AI, graph neural networks, and HPC simulation.}}
\vspace{1pt}
\item{\cventry{Aug 2016--Jul 2019}{M.S., Computer and Information Sciences}{Delaware State University}{Dover, DE}{}{GPA: 3.9; machine learning, pattern recognition, and HPC.}}
\vspace{1pt}
\item{\cventry{2003--2005}{M.E., Software Engineering}{Tsinghua University}{Beijing, China}{}{Software architecture, enterprise systems, and project management.}}
\vspace{1pt}
\item{\cventry{1998--2002}{B.E., Ocean Engineering; Minor in Computer Science}{Dalian University of Technology}{Dalian, China}{}{Engineering, applied mathematics, and computer science.}}
\end{itemize}

\section{Selected Publications}
\vspace{0pt}
\begin{itemize}
\item J.~Zhao, K.~Kong, X.~Hei, Y.~Tu, X.~Du, ``A Visible Light Channel based Access Control Scheme for Wireless Insulin Pump Systems,'' \textit{Proc.\ IEEE ICC}, 2018.
\item J.~Zhao, L.~Lu, Z.~Tyler, K.~Kong, Q.~Lu, J.~Ren, ``Simulated Raman Spectral Analysis of Organic Molecules,'' \textit{Proc.\ SPIE} 10526, 105262I, 2018. doi:~10.1117/12.2291176
\item A.~Bhat, J.~Zhao, T.~Cooks, J.~Ren, Q.~Lu, ``Dark-field Microscopic Study of the Interactions between Gold/Silver Nanoparticles and Giant Unilamellar Vesicles,'' \textit{Proc.\ SPIE} 10506, 105061C, 2018. doi:~10.1117/12.2290828
\item Co-translator, W.~van~der~Aalst \& K.~van~Hee, \textit{Workflow Management: Models, Methods, and Systems} (Chinese ed.), Tsinghua Univ.\ Press, 2004. (Supervised by Prof.\ Jianmin Wang; co-translated with L.~Wen et al.)
\item Additional papers in workflow management (colored Petri nets, Tsinghua 2004--2005) and energy-pipeline analytics (PetroChina 2007--2016); full list available on request.
\end{itemize}

\section{Additional}
\vspace{0pt}
\begin{itemize}
\item \textbf{Languages}: Chinese/Mandarin (native); English (professional working proficiency).
\item \textbf{Work Authorization}: H-1B transfer available.
\item \textbf{Honors \& Certifications}:
  \begin{itemize}
    \item FIRST LEGO League (FLL) Certified Judge: Judge Training and Robot Design Judge Training (March 2020)
    \item CITI Program, Responsible Conduct of Research for Engineers, Delaware State University (Jan.\ 2018)
    \item Vice President, Science \& Technology Association, Tsinghua University School of Software (2003--2005)
    \item Top 20\% academic ranking, M.E.\ program, Tsinghua University
  \end{itemize}
\item \textbf{References}: Prof.\ Jun Ren, Delaware State University (\href{mailto:jren@desu.edu}{jren@desu.edu}); Dr.\ Xiali (Sharon) Hei, Alfred and Helen M.\ Lamson Endowed Professor in Computer Science, University of Louisiana at Lafayette (\href{mailto:xiali.hei@louisiana.edu}{xiali.hei@louisiana.edu}; +1-337-482-1037). Additional references available on request.
\end{itemize}

\end{document}
